\section{Shor's Algorithm: Impact on Blockchain Primitives}

\label{sec:shor}

\subsection{Algorithm Overview}

Shor's algorithm~\cite{shor1997} finds the order of an element in a finite group
using quantum phase estimation. For an element $g$ of order $r$ in a group
$\mathbb{G}$, Shor's algorithm finds $r$ in $O((\log |\mathbb{G}|)^2 \cdot \log\log |\mathbb{G}|)$
quantum gates, with $O(\log |\mathbb{G}|)$ qubits. Given the order, the discrete
logarithm follows by standard reductions.

\subsection{ECDSA (secp256k1)}

ECDSA security rests on the difficulty of computing $k$ from $kG$ where $G$ is the
generator of a 256-bit elliptic curve group. Shor's algorithm solves this with
approximately $2{,}330$ logical qubits and $1.26 \times 10^{11}$ T-gates for
secp256k1~\cite{roetteler2017quantum}. With a $1\,\mu$s T-gate time and a
surface code implementing error correction:

\begin{table}[ht]
\centering
\caption{Quantum resource estimates for breaking elliptic curve DLP.}
\label{tab:ecdsa-resources}
\begin{tabular}{lrrr}
\toprule
\textbf{Curve} & \textbf{Logical Qubits} & \textbf{T-gates} & \textbf{Est. Time} \\
\midrule
secp256k1 (256-bit) & 2{,}330 & $1.26 \times 10^{11}$ & 8 hours \\
Ed25519 (255-bit) & 2{,}310 & $1.24 \times 10^{11}$ & 8 hours \\
P-384 (384-bit) & 3{,}484 & $4.16 \times 10^{11}$ & 26 hours \\
\bottomrule
\end{tabular}
\vspace{0.3em}
\footnotesize{Resource estimates from Roetteler et al.~\cite{roetteler2017quantum},
assuming surface code with code distance 27 and $1\,\mu$s cycle time.}
\end{table}

The critical observation: once a quantum computer exists with sufficient qubits, key
recovery is measured in \emph{hours}, not years. There is no gradual degradation.

\subsection{BLS12-381}

BLS signatures operate on the pairing-friendly curve BLS12-381, where security rests
on the co-CDH assumption in groups $\mathbb{G}_1$ (381-bit) and $\mathbb{G}_2$
(762-bit). Shor's algorithm applies to both groups. Breaking BLS12-381 requires
approximately $3{,}800$ logical qubits---more than secp256k1 due to the larger
field, but the same order of magnitude.

BLS is more dangerous to break than ECDSA in the blockchain context because:
\begin{itemize}[leftmargin=2em]
  \item Validator BLS keys are \emph{static}: they persist for the validator's
        entire lifetime, often years.
  \item BLS aggregate signatures compress $n$ validator votes into 48 bytes. If BLS
        is broken, an attacker can forge consensus certificates for any historical or
        future block.
  \item A forged BLS certificate can rewrite chain history or create conflicting
        finalized chains---a safety violation, not merely a liveness issue.
\end{itemize}

\subsection{RSA}

While RSA is not commonly used in blockchain consensus, it appears in:
\begin{itemize}[leftmargin=2em]
  \item TLS certificates for RPC endpoints
  \item Verifiable delay functions (VDFs) based on RSA groups
  \item Some oracle attestation schemes
\end{itemize}

Shor's algorithm factors RSA-2048 with approximately $4{,}098$ logical qubits.
RSA-3072 requires approximately $6{,}146$ logical qubits. Both are within the
projected capability of CRQCs by 2035.

%% ========================================================================
%%  3. GROVER'S ALGORITHM: HASH FUNCTION IMPACT
%% ========================================================================
